\begin{theorem}[Boundary rigidity]\label{thm:main}
Let $n\ge2$ be an integer. Every facet of $R_n(c)$ is isometric to $R_{n-1}(c)$ when $n\ge3$, and to a unit segment when $n=2$.
If $n\ge3$, the intrinsic piecewise flat metric of the boundary of $R_n(c)$ determines $c$, and hence determines $R_n(c)$ up to congruence. If $n=2$, the intrinsic metric of the boundary is that of a circle of length $4$ for every $c$, so $c$ is not determined, although the area $\sqrt{1-c^2}$ varies with $c$.
\end{theorem}
\begin{proof}
The facet statement is Lemma~\ref{lem:facet}.
Let $n\ge3$. Two edges at a vertex lie in a common flat $2$-face, so the intrinsic metric of the boundary determines the angle between any two edges at any vertex, hence the Gram matrices $D_\varepsilon G_n(c)D_\varepsilon$ of Lemma~\ref{lem:vertex} at all vertices, up to permutation of the edges. The vertex $S=\emptyset$ has all off-diagonal entries equal to $c$. By Lemma~\ref{lem:vertex}, any vertex whose off-diagonal entries are all equal has that common value $c$, so $c$ is read off from the boundary. The pair $(n,c)$ determines $R_n(c)$ up to an orthogonal map by the Cholesky argument: if $W^\T W=V^\T V$ with $V$ invertible then $W=QV$ with $Q=WV^{-1}$ orthogonal. For $n=2$ no such recovery is possible, as the boundary is a circle of length $4$ whatever $c$ is.
\end{proof}
